\input{preamble}

\title{Section 1.8 --- Absolute Value Equations --- Answer Key}

\begin{document}
\maketitle

%\section*{Evaluating Absolute Values}
%Evaluate each expression.
%\begin{smenum}
%\mitemxx
%{$\Abs{2-3}+\dfrac{\Abs{5}}{\Abs{-4}}=\frac{9}{4}$}
%{$-\Abs{-5}\cdot\Abs{7-1}=-30$}
%\mitemxx
%{$2\cdot-\Abs{6+3}-\Abs{-2\cdot 6}=-30$}
%{$-\frac{1}{3}\Abs{2-8}+4=2$}
%\end{smenum}
%
%%%%%%%%%%%%%%%%%%%%%%%%%%%%%%%%%%%%%%%%%%%%%%%%%%%%%%%%%
%\section*{Identifying Absolute Value Functions}
%Determine whether each of the following is an absolute value function.
%\begin{smenum}
%\mitemxxxx
%{Yes}
%{No}
%{No}
%{Yes}
%\mitemxxxx
%{No}
%{No}
%{Yes}
%{No}
%\end{smenum}
%
%%%%%%%%%%%%%%%%%%%%%%%%%%%%%%%%%%%%%%%%%%%%%%%%%%%%%%%%%
%\section*{Graphing Absolute Value Functions}
%Graph each absolute value function. Then, looking at the graph, state the domain and range of each function.
%\begin{smenum}
%\mitemxx
%{$f(x)=\frac{2}{3}\Abs{x}$\\
%$\dom f=\R$\quad$\rng f=\Lint{0}$\\
%\begin{ilpic}[x=0.2in,y=0.2in]
%\DrawGridSquare
%\draw[<->,thick] (-5.9,2/3*5.9) -- (0,0) -- (5.9,2/3*5.9);
%\foreach \x/\y in {-3/2,0/0,3/2} \fill (\x,\y) circle (0.15);
%\end{ilpic}}
%{$g(x)=2\Abs{x-4}-2$\\
%$\dom g=\R$\quad$\rng g=\Lint{-2}$\\
%\begin{ilpic}[x=0.2in,y=0.2in]
%%\useasboundingbox(-5.9,-5.9) rectangle (5.9,5.9);
%\DrawGridSquare
%\draw[<->,thick] (0.05,5.9) -- (4,-2) -- (5.9,1.8);
%\foreach \x/\y in {3/0,4/-2,5/0} \fill (\x,\y) circle (0.15);
%\end{ilpic}}
%\mitemxx
%{$h(x)=-\frac{3}{4}\Abs{x}+1$\\
%$\dom h=\R$\quad$\rng h=\Rint{1}$\\
%\begin{ilpic}[x=0.2in,y=0.2in]
%%\useasboundingbox(-5.9,-5.9) rectangle (5.9,5.9);
%\DrawGridSquare
%\draw[<->,thick] (-5.9,-3.425) -- (0,1) -- (5.9,-3.425);
%\foreach \x/\y in {-4/-2,0/1,4/-2} \fill (\x,\y) circle (0.15);
%\end{ilpic}}
%{$k(x)=-\frac{3}{2}\Abs{x+1}-1$\\
%$\dom k=\R$\quad$\rng k=\Rint{-1}$\\
%\begin{ilpic}[x=0.2in,y=0.2in]
%%\useasboundingbox(-5.9,-5.9) rectangle (5.9,5.9);
%\DrawGridSquare
%\draw[<->,thick] (-12.8/3,-5.9) -- (-1,-1) -- (6.8/3,-5.9);
%\foreach \x/\y in {-3/-4,-1/-1,1/-4} \fill (\x,\y) circle (0.15);
%\end{ilpic}}
%\end{smenum}
%
%%%%%%%%%%%%%%%%%%%%%%%%%%%%%%%%%%%%%%%%%%%%%%%%%%%%%%%%%
%\section*{Range and Domain of an absolute value function}
%Find the domain and range of each function.
%\begin{smenum}
%\mitemxxxx
%{$\dom f=\R$\\$\rng f=\Lint{0}$}
%{$\dom g=\R$\\$\rng g=\Lint{4}$}
%{$\dom h=\R$\\$\rng h=\Rint{3}$}
%{$\dom k=\R$\\$\rng k=\Lint{-4}$}
%\end{smenum}

%%%%%%%%%%%%%%%%%%%%%%%%%%%%%%%%%%%%%%%%%%%%%%%%%%%%%%%%
\section*{Zero, One, or Two Solutions}
Determine whether each absolute value equation has zero, one, or two solutions.  If the equation has one or two solutions, find them.
\begin{smenum}
\mitemxxxx
{Two solutions\\$x=-4$ or $x=4$}
{Zero solutions}
{One solution\\$x=-6$}
{Two solutions\\$x=-11$ or $x=23$}
\mitemxxxx
{Zero solutions}
{Two solutions\\$x=-\frac{27}{2}$ or $x=\frac{3}{2}$}
{One solution\\$x=\frac{1}{4}$}
{Zero solutions}
\end{smenum}

%%%%%%%%%%%%%%%%%%%%%%%%%%%%%%%%%%%%%%%%%%%%%%%%%%%%%%%%
\section*{Solve Absoulte Value Equations}
Find any possible solutions to each equation.
\begin{smenum}
\mitemxxxx
{$r=-1$ or $r=1$}
{$x=-5$ or $x=-1$}
{No solution}
{No solution}
\mitemxxxx
{$x=0$ or $x=10$}
{$y=-1$ or $y=5$}
{$u=-9$ or $u=15$}
{$p=-\frac{11}{3}$ or $p=\frac{7}{3}$}
\mitemxx
{$x=-\frac{3}{2}$}
{$x=6$ or $x=48$}
\end{smenum}

\end{document}